%! Functions and Function Notation — Unit 1, Lesson 1.1 — Warm-Up
% Exactly ONE page, blank and keyed (LESSON_SHAPE.md, one_page). Four prerequisite items:
% 1–3 are the algebra today's evaluating runs on; item 4 is the inequality-membership skill
% every piecewise condition depends on.
\documentclass[10pt]{article}
\usepackage{precalculus-article}
\usepackage{precalculus-boxes}
\newcommand{\TallMath}[1]{$\displaystyle #1\rule[-1.4em]{0pt}{3.2em}$}

\begin{document}

\pageheader{Unit 1, Lesson 1.1}{Warm-Up}

\vspace{0.12in}

\begin{enumerate}[itemsep=10pt]

\item Evaluate $4x - 9$ when $x = -3$.

\vspace{0.1in}
Value: \blank{3cm}

\item Evaluate $x^2 + 2x$ when $x = -5$.

\begin{work}
  (-5)^2 + 2(-5) &= 25 - 10 \\
                 &= 15
\end{work}

Value: \blank{3cm}

\item Solve for $x$: \quad $3x - 5 = 16$

\begin{work}
  3x - 5 &= 16 \\
      3x &= 21 \\
       x &= 7
\end{work}

$x = $ \blank{3cm}

\item A condition like $x \ge 2$ is a test each number either passes or fails.
Circle every number that \textbf{passes}.

\begin{enumerate}[label=(\alph*), itemsep=6pt]
  \item $x \ge 2$: \qquad $-1$ \qquad $0$ \qquad $2$ \qquad $5$
  \item $x < 2$: \qquad $-1$ \qquad $0$ \qquad $2$ \qquad $5$
\end{enumerate}

\end{enumerate}

\end{document}
